This paper studies local scene emergence in heavy metal. The question is not why some cities have
more bands in general. It is why some city-genre environments move from scattered local entry to a
recognizable local scene earlier than others. I treat that process as a localized production
problem in a project-based creative industry, where bands are projects, musicians move across
teams, and new local niches are easier to sustain once enough specialized capability has
accumulated.

That question is difficult to study with the usual retrospective approach. Famous scenes are often
identified only after genre labels are settled, canonical bands are known, and local histories can
be written in hindsight. But genre recognition is not the same thing as genre formation. By the
time a label is stabilized, much of the relevant local infrastructure may already have formed
through repeated collaboration, spinouts, and worker mobility across projects. The empirical
challenge is therefore to measure local scene formation before the historical narrative fully
settles.

Heavy metal offers an unusually useful setting for that task. Metal Archives provides a broad band
census with genre strings, formation years, and location information, while the member dump makes
it possible to observe overlapping personnel ties across bands. I use those sources to build a
city-level collaboration dataset from 971,877 musician-band edges covering 682,126 musicians, and
I link that object to a dynamic city-by-genre panel of local scene emergence. The resulting panel
contains 56,978 city-genre cells, of which 5,756 eventually cross the paper's operational
emergence threshold. The preferred exact-year fixed-effects sample contains 171,653
city-genre-years and 5,306 emergence events.

The paper's outcome is intentionally operational. A city-genre cell is counted as emergent in the
year it first crosses the paper's fifth-band threshold in that broad genre family. This is not a
claim that genre identity literally begins at band number five. It is a transparent empirical rule
for the transition from isolated local entry to operational scene status. That transparency is a
feature of the design, but it also defines the paper's main empirical risk: one must show that the
results are not merely a mechanical consequence of thicker places crossing a stock threshold
sooner.

Two empirical facts organize the paper. The first is descriptive. Exact-year emergence rates rise
sharply with local band stock and with the depth of the local multi-band musician pool. The second
is reduced-form. In exact-year fixed-effects transition specifications, later city-genre
transitions into operational scene status track three stable lagged predictors: local spawning
flow, target-genre active bands, and target-genre multi-band musician depth. Composition-heavy
extensions and broader connector measures appear only in background work because they do not
improve the main paper object. Threshold-response checks show that this pattern is not pinned to
the fifth-band cutoff, but they also narrow the claim: the paper is strongest on the later
approach to emergence rather than on the earliest seed stage of scene formation.

The contribution is therefore narrower than the broadest version of the project, but clearer.
First, the paper introduces a large city-level collaboration dataset for metal built from the full
band census and member records rather than from a curated set of famous scenes. Second, it treats
genre emergence as a measurable local event that can be studied in panel form before retrospective
scene histories dominate the evidence. Third, it shows in that panel that later transitions into
operational scene status are most closely associated with local spawning, focal-niche band stock,
and focal-niche overlapping-worker depth, with the clearest evidence coming near the later
approach to the scene threshold.

The paper remains reduced-form rather than causal. It does not identify an exogenous source of
scene formation, observe knowledge transfer directly, or date the true historical origin of a
genre. The contribution is instead a carefully validated empirical pattern. In a large panel of
not-yet-emerged city-genre cells, transitions into operational scene status are more likely where
local capability is thicker, project spinouts are more common, and workers are already recombining
within related local projects. Heavy metal is useful here not because it is culturally
exceptional, but because the setting makes project overlap and worker mobility unusually
observable.

The rest of the paper proceeds as follows. Section 2 situates the paper in the related
literature. Section 3 presents the empirical motivation. Section 4 sets out a compact model of
local niche formation. Section 5 describes the data and sample construction. Section 6 presents the
empirical approach. Section 7 reports the descriptive and reduced-form results. The appendix
collects the detailed sample-construction and variable-definition objects behind the main
specification.
